\section{Metalinguistic Feature Integration}
\label{sec:metalinguistic-analysis}

The continued pre-training experiments in Section~\ref{sec:continued-pretraining-results} establish that metalinguistic auxiliary objectives can improve article-level ideology classification, though only under specific configurations. Rather than re-stating the per-configuration F1 numbers, we focus here on what those findings imply about the role of metalinguistic structure in ideological representation.

Three patterns merit particular attention. Theme prediction is the auxiliary objective that delivers the largest and most consistent gains, with the clearest improvements on the historically difficult center and right categories. This is consistent with the cognitive-science framing that topical emphasis is a structural feature of how outlets encode stance, and it suggests that theme complements, rather than substitutes for, the lexical content captured by contrastive pre-training. The benefit of auxiliary objectives is, however, contingent on a sufficiently dense primary contrastive signal: with sparse triplet supervision the auxiliary heads inject more noise than signal. Multi-task pre-training of this form is therefore best understood not as a remedy for weak contrastive supervision but as an amplifier of a sufficiently strong one. Finally, the gains do not stack additively. Combining theme and tone often underperforms theme alone, indicating that the two metalinguistic dimensions overlap in the discriminative information they contribute and that naive summation of auxiliary losses is not the appropriate integration strategy.

Taken together, these results validate the central premise of Chapter~\ref{sec:continued-pretraining}: metalinguistic structure carries ideological signal that is not fully captured by a content-only contrastive objective, and auxiliary metalinguistic supervision is worth including in future continued pre-training work, provided that the primary objective and its supervision density are chosen with care. The gains measured in absolute F1 terms are modest but consistent. Whether they translate into corresponding improvements on the evaluation axes that F1 does not directly measure, namely robustness to shortcut learning and alignment with the human causal profile, is the question the remainder of this chapter takes up.
